% paper99-copilot-connection-health.tex — Copilot Connection Hooks and Health Checks
%   CopilotConnectionHook + CopilotHealthCheck
%   Paper 99 of the Judgment Geometry series.
% Compile: pdflatex paper99-copilot-connection-health
\documentclass[11pt]{article}
\usepackage{lmodern}
% jugeo-common.tex — shared preamble for all Judgment Geometry papers
% Include via: % jugeo-common.tex — shared preamble for all Judgment Geometry papers
% Include via: % jugeo-common.tex — shared preamble for all Judgment Geometry papers
% Include via: \input{jugeo-common}

% ── Packages ────────────────────────────────────────────────────────────
\usepackage{amsmath,amssymb,amsthm,mathtools}
\usepackage{tikz-cd}
\usepackage{listings}
\usepackage{xcolor}
\usepackage{xspace}
\usepackage{booktabs}
\usepackage{multirow}
\usepackage{enumitem}
\usepackage[expansion=false]{microtype}
\usepackage{hyperref}
\usepackage{cleveref}
% \usepackage{thmtools}  % disabled: not available in this TeX installation
\usepackage{float}
\usepackage{caption}
\usepackage{subcaption}
\usepackage{tabularx}
\usepackage{stmaryrd}       % \llbracket, \rrbracket
\usepackage{bbm}            % \mathbbm{1}

% ── Page geometry ───────────────────────────────────────────────────────
\usepackage[margin=1in]{geometry}

% ── Theorem environments ────────────────────────────────────────────────
\theoremstyle{plain}
\newtheorem{theorem}{Theorem}[section]
\newtheorem{lemma}[theorem]{Lemma}
\newtheorem{proposition}[theorem]{Proposition}
\newtheorem{corollary}[theorem]{Corollary}

\theoremstyle{definition}
\newtheorem{definition}[theorem]{Definition}
\newtheorem{example}[theorem]{Example}
\newtheorem{construction}[theorem]{Construction}

\theoremstyle{remark}
\newtheorem{remark}[theorem]{Remark}
\newtheorem{notation}[theorem]{Notation}

% ── Python listings style ──────────────────────────────────────────────
\definecolor{jugeoBlue}{HTML}{2563EB}
\definecolor{jugeoGreen}{HTML}{059669}
\definecolor{jugeoRed}{HTML}{DC2626}
\definecolor{jugeoPurple}{HTML}{7C3AED}
\definecolor{jugeoGray}{HTML}{6B7280}
\definecolor{jugeoBg}{HTML}{F8FAFC}

\lstdefinestyle{jugeo-python}{
  language=Python,
  basicstyle=\ttfamily\small,
  keywordstyle=\color{jugeoBlue}\bfseries,
  stringstyle=\color{jugeoGreen},
  commentstyle=\color{jugeoGray}\itshape,
  numberstyle=\tiny\color{jugeoGray},
  backgroundcolor=\color{jugeoBg},
  numbers=left,
  numbersep=6pt,
  frame=single,
  framerule=0.4pt,
  rulecolor=\color{jugeoGray!40},
  breaklines=true,
  breakatwhitespace=true,
  showstringspaces=false,
  tabsize=4,
  xleftmargin=1.5em,
  framexleftmargin=1.5em,
  morekeywords={dataclass,frozen,field,Enum,IntEnum,auto,Any,
                Mapping,Sequence,Optional,match,case},
  literate={→}{{$\to$}}1 {←}{{$\leftarrow$}}1
           {≤}{{$\leq$}}1 {≥}{{$\geq$}}1 {≠}{{$\neq$}}1
           {∈}{{$\in$}}1 {∀}{{$\forall$}}1 {∃}{{$\exists$}}1
           {⊕}{{$\oplus$}}1 {⊖}{{$\ominus$}}1,
  escapeinside={(*@}{@*)},
}
\lstset{style=jugeo-python}

\lstdefinestyle{jugeo-output}{
  basicstyle=\ttfamily\small,
  backgroundcolor=\color{jugeoBg},
  frame=single,
  framerule=0.4pt,
  rulecolor=\color{jugeoGray!40},
  breaklines=true,
  showstringspaces=false,
  xleftmargin=1.5em,
  framexleftmargin=0.5em,
  numbers=none,
}

% ── JuGeo-specific macros ──────────────────────────────────────────────

% Judgment 8-tuple
\newcommand{\J}{\mathcal{J}}
\newcommand{\Jud}[8]{(#1,\,#2,\,#3,\,#4,\,#5,\,#6,\,#7,\,#8)}
\newcommand{\JudShort}[1]{\J_{#1}}

% Components
\newcommand{\coord}{\mathsf{c}}            % coordinate
\newcommand{\prop}{\varphi}                 % proposition
\newcommand{\carrier}{\mathcal{A}}          % carrier type
\newcommand{\evid}{\mathcal{E}}             % evidence bundle
\newcommand{\oblig}{\mathcal{O}}            % obligations
\newcommand{\obstr}{\mathcal{B}}            % obstructions
\newcommand{\trust}{\mathcal{T}}            % trust annotation
\newcommand{\prov}{\Pi}                     % provenance

% Trust algebra
\newcommand{\Talg}{\mathfrak{T}}            % trust ordered algebra
\newcommand{\Eadm}{\mathcal{E}_{\mathrm{adm}}}
\newcommand{\tleq}{\preceq}                 % trust ordering
\newcommand{\tjoin}{\oplus}                 % conservative join
\newcommand{\tatten}{\ominus}               % attenuation
\newcommand{\tpromote}{\uparrow_\pi}        % promotion
\newcommand{\tdemote}{\downarrow_\chi}       % demotion

% Trust tiers
\newcommand{\Tcontra}{\ensuremath{\mathsf{CONTRADICTED}}}
\newcommand{\Tunver}{\ensuremath{\mathsf{UNVERIFIED}}}
\newcommand{\Tcopilot}{\ensuremath{\mathsf{COPILOT}}}
\newcommand{\Toracle}{\ensuremath{\mathsf{ORACLE}}}
\newcommand{\Truntime}{\ensuremath{\mathsf{RUNTIME}}}
\newcommand{\Tsolver}{\ensuremath{\mathsf{SOLVER}}}
\newcommand{\Tproof}{\ensuremath{\mathsf{PROOF}}}

% Site and geometry
\newcommand{\Site}{\mathcal{S}}             % semantic site
\newcommand{\Cov}{\mathrm{Cov}}            % covering family
\newcommand{\Ob}{\mathrm{Ob}}              % objects
\newcommand{\Mor}{\mathrm{Mor}}            % morphisms
\newcommand{\res}{\mathrm{res}}            % restriction
\newcommand{\Cech}{\check{C}}              % Čech complex
\newcommand{\Hcoh}[1]{H^{#1}}             % cohomology group

% Descent
\newcommand{\descent}{\mathrm{desc}}
\newcommand{\glue}{\mathrm{glue}}
\newcommand{\overlap}[2]{#1 \cap #2}

% Semantic moves
\newcommand{\VERIFY}{\textsc{Verify}}
\newcommand{\CONSTRUCT}{\textsc{Construct}}
\newcommand{\NORMALIZE}{\textsc{Normalize}}
\newcommand{\GLUE}{\textsc{Glue}}
\newcommand{\RESTRICT}{\textsc{Restrict}}
\newcommand{\TRANSPORT}{\textsc{Transport}}
\newcommand{\REPAIR}{\textsc{Repair}}
\newcommand{\PROMOTE}{\textsc{Promote}}
\newcommand{\CHALLENGE}{\textsc{Challenge}}
\newcommand{\ESCALATE}{\textsc{Escalate}}

% Fragments
\newcommand{\QFLIA}{\texttt{QF\_LIA}}
\newcommand{\QFLRA}{\texttt{QF\_LRA}}
\newcommand{\QFBV}{\texttt{QF\_BV}}
\newcommand{\QFUF}{\texttt{QF\_UF}}

% Misc
\newcommand{\Python}{\textsc{Python}}
\newcommand{\JuGeo}{\textsc{JuGeo}}
\newcommand{\LEAN}{\textsc{Lean}\,4}
\newcommand{\Fstar}{F$^{\star}$}
\newcommand{\Coq}{\textsc{Coq}}
\newcommand{\Dafny}{\textsc{Dafny}}

% Common shortcuts
\newcommand{\ie}{i.e.\@\xspace}
\newcommand{\eg}{e.g.\@\xspace}
\newcommand{\cf}{cf.\@\xspace}
\newcommand{\wrt}{w.r.t.\@\xspace}

% Evaluation shortcuts
\newcommand{\numcases}{300}
\newcommand{\numspec}{100}
\newcommand{\numeq}{100}
\newcommand{\numbug}{100}
\newcommand{\medtime}{6.5\,ms}
\newcommand{\totaltime}{1.949\,s}


% ── Packages ────────────────────────────────────────────────────────────
\usepackage{amsmath,amssymb,amsthm,mathtools}
\usepackage{tikz-cd}
\usepackage{listings}
\usepackage{xcolor}
\usepackage{xspace}
\usepackage{booktabs}
\usepackage{multirow}
\usepackage{enumitem}
\usepackage[expansion=false]{microtype}
\usepackage{hyperref}
\usepackage{cleveref}
% \usepackage{thmtools}  % disabled: not available in this TeX installation
\usepackage{float}
\usepackage{caption}
\usepackage{subcaption}
\usepackage{tabularx}
\usepackage{stmaryrd}       % \llbracket, \rrbracket
\usepackage{bbm}            % \mathbbm{1}

% ── Page geometry ───────────────────────────────────────────────────────
\usepackage[margin=1in]{geometry}

% ── Theorem environments ────────────────────────────────────────────────
\theoremstyle{plain}
\newtheorem{theorem}{Theorem}[section]
\newtheorem{lemma}[theorem]{Lemma}
\newtheorem{proposition}[theorem]{Proposition}
\newtheorem{corollary}[theorem]{Corollary}

\theoremstyle{definition}
\newtheorem{definition}[theorem]{Definition}
\newtheorem{example}[theorem]{Example}
\newtheorem{construction}[theorem]{Construction}

\theoremstyle{remark}
\newtheorem{remark}[theorem]{Remark}
\newtheorem{notation}[theorem]{Notation}

% ── Python listings style ──────────────────────────────────────────────
\definecolor{jugeoBlue}{HTML}{2563EB}
\definecolor{jugeoGreen}{HTML}{059669}
\definecolor{jugeoRed}{HTML}{DC2626}
\definecolor{jugeoPurple}{HTML}{7C3AED}
\definecolor{jugeoGray}{HTML}{6B7280}
\definecolor{jugeoBg}{HTML}{F8FAFC}

\lstdefinestyle{jugeo-python}{
  language=Python,
  basicstyle=\ttfamily\small,
  keywordstyle=\color{jugeoBlue}\bfseries,
  stringstyle=\color{jugeoGreen},
  commentstyle=\color{jugeoGray}\itshape,
  numberstyle=\tiny\color{jugeoGray},
  backgroundcolor=\color{jugeoBg},
  numbers=left,
  numbersep=6pt,
  frame=single,
  framerule=0.4pt,
  rulecolor=\color{jugeoGray!40},
  breaklines=true,
  breakatwhitespace=true,
  showstringspaces=false,
  tabsize=4,
  xleftmargin=1.5em,
  framexleftmargin=1.5em,
  morekeywords={dataclass,frozen,field,Enum,IntEnum,auto,Any,
                Mapping,Sequence,Optional,match,case},
  literate={→}{{$\to$}}1 {←}{{$\leftarrow$}}1
           {≤}{{$\leq$}}1 {≥}{{$\geq$}}1 {≠}{{$\neq$}}1
           {∈}{{$\in$}}1 {∀}{{$\forall$}}1 {∃}{{$\exists$}}1
           {⊕}{{$\oplus$}}1 {⊖}{{$\ominus$}}1,
  escapeinside={(*@}{@*)},
}
\lstset{style=jugeo-python}

\lstdefinestyle{jugeo-output}{
  basicstyle=\ttfamily\small,
  backgroundcolor=\color{jugeoBg},
  frame=single,
  framerule=0.4pt,
  rulecolor=\color{jugeoGray!40},
  breaklines=true,
  showstringspaces=false,
  xleftmargin=1.5em,
  framexleftmargin=0.5em,
  numbers=none,
}

% ── JuGeo-specific macros ──────────────────────────────────────────────

% Judgment 8-tuple
\newcommand{\J}{\mathcal{J}}
\newcommand{\Jud}[8]{(#1,\,#2,\,#3,\,#4,\,#5,\,#6,\,#7,\,#8)}
\newcommand{\JudShort}[1]{\J_{#1}}

% Components
\newcommand{\coord}{\mathsf{c}}            % coordinate
\newcommand{\prop}{\varphi}                 % proposition
\newcommand{\carrier}{\mathcal{A}}          % carrier type
\newcommand{\evid}{\mathcal{E}}             % evidence bundle
\newcommand{\oblig}{\mathcal{O}}            % obligations
\newcommand{\obstr}{\mathcal{B}}            % obstructions
\newcommand{\trust}{\mathcal{T}}            % trust annotation
\newcommand{\prov}{\Pi}                     % provenance

% Trust algebra
\newcommand{\Talg}{\mathfrak{T}}            % trust ordered algebra
\newcommand{\Eadm}{\mathcal{E}_{\mathrm{adm}}}
\newcommand{\tleq}{\preceq}                 % trust ordering
\newcommand{\tjoin}{\oplus}                 % conservative join
\newcommand{\tatten}{\ominus}               % attenuation
\newcommand{\tpromote}{\uparrow_\pi}        % promotion
\newcommand{\tdemote}{\downarrow_\chi}       % demotion

% Trust tiers
\newcommand{\Tcontra}{\ensuremath{\mathsf{CONTRADICTED}}}
\newcommand{\Tunver}{\ensuremath{\mathsf{UNVERIFIED}}}
\newcommand{\Tcopilot}{\ensuremath{\mathsf{COPILOT}}}
\newcommand{\Toracle}{\ensuremath{\mathsf{ORACLE}}}
\newcommand{\Truntime}{\ensuremath{\mathsf{RUNTIME}}}
\newcommand{\Tsolver}{\ensuremath{\mathsf{SOLVER}}}
\newcommand{\Tproof}{\ensuremath{\mathsf{PROOF}}}

% Site and geometry
\newcommand{\Site}{\mathcal{S}}             % semantic site
\newcommand{\Cov}{\mathrm{Cov}}            % covering family
\newcommand{\Ob}{\mathrm{Ob}}              % objects
\newcommand{\Mor}{\mathrm{Mor}}            % morphisms
\newcommand{\res}{\mathrm{res}}            % restriction
\newcommand{\Cech}{\check{C}}              % Čech complex
\newcommand{\Hcoh}[1]{H^{#1}}             % cohomology group

% Descent
\newcommand{\descent}{\mathrm{desc}}
\newcommand{\glue}{\mathrm{glue}}
\newcommand{\overlap}[2]{#1 \cap #2}

% Semantic moves
\newcommand{\VERIFY}{\textsc{Verify}}
\newcommand{\CONSTRUCT}{\textsc{Construct}}
\newcommand{\NORMALIZE}{\textsc{Normalize}}
\newcommand{\GLUE}{\textsc{Glue}}
\newcommand{\RESTRICT}{\textsc{Restrict}}
\newcommand{\TRANSPORT}{\textsc{Transport}}
\newcommand{\REPAIR}{\textsc{Repair}}
\newcommand{\PROMOTE}{\textsc{Promote}}
\newcommand{\CHALLENGE}{\textsc{Challenge}}
\newcommand{\ESCALATE}{\textsc{Escalate}}

% Fragments
\newcommand{\QFLIA}{\texttt{QF\_LIA}}
\newcommand{\QFLRA}{\texttt{QF\_LRA}}
\newcommand{\QFBV}{\texttt{QF\_BV}}
\newcommand{\QFUF}{\texttt{QF\_UF}}

% Misc
\newcommand{\Python}{\textsc{Python}}
\newcommand{\JuGeo}{\textsc{JuGeo}}
\newcommand{\LEAN}{\textsc{Lean}\,4}
\newcommand{\Fstar}{F$^{\star}$}
\newcommand{\Coq}{\textsc{Coq}}
\newcommand{\Dafny}{\textsc{Dafny}}

% Common shortcuts
\newcommand{\ie}{i.e.\@\xspace}
\newcommand{\eg}{e.g.\@\xspace}
\newcommand{\cf}{cf.\@\xspace}
\newcommand{\wrt}{w.r.t.\@\xspace}

% Evaluation shortcuts
\newcommand{\numcases}{300}
\newcommand{\numspec}{100}
\newcommand{\numeq}{100}
\newcommand{\numbug}{100}
\newcommand{\medtime}{6.5\,ms}
\newcommand{\totaltime}{1.949\,s}


% ── Packages ────────────────────────────────────────────────────────────
\usepackage{amsmath,amssymb,amsthm,mathtools}
\usepackage{tikz-cd}
\usepackage{listings}
\usepackage{xcolor}
\usepackage{xspace}
\usepackage{booktabs}
\usepackage{multirow}
\usepackage{enumitem}
\usepackage[expansion=false]{microtype}
\usepackage{hyperref}
\usepackage{cleveref}
% \usepackage{thmtools}  % disabled: not available in this TeX installation
\usepackage{float}
\usepackage{caption}
\usepackage{subcaption}
\usepackage{tabularx}
\usepackage{stmaryrd}       % \llbracket, \rrbracket
\usepackage{bbm}            % \mathbbm{1}

% ── Page geometry ───────────────────────────────────────────────────────
\usepackage[margin=1in]{geometry}

% ── Theorem environments ────────────────────────────────────────────────
\theoremstyle{plain}
\newtheorem{theorem}{Theorem}[section]
\newtheorem{lemma}[theorem]{Lemma}
\newtheorem{proposition}[theorem]{Proposition}
\newtheorem{corollary}[theorem]{Corollary}

\theoremstyle{definition}
\newtheorem{definition}[theorem]{Definition}
\newtheorem{example}[theorem]{Example}
\newtheorem{construction}[theorem]{Construction}

\theoremstyle{remark}
\newtheorem{remark}[theorem]{Remark}
\newtheorem{notation}[theorem]{Notation}

% ── Python listings style ──────────────────────────────────────────────
\definecolor{jugeoBlue}{HTML}{2563EB}
\definecolor{jugeoGreen}{HTML}{059669}
\definecolor{jugeoRed}{HTML}{DC2626}
\definecolor{jugeoPurple}{HTML}{7C3AED}
\definecolor{jugeoGray}{HTML}{6B7280}
\definecolor{jugeoBg}{HTML}{F8FAFC}

\lstdefinestyle{jugeo-python}{
  language=Python,
  basicstyle=\ttfamily\small,
  keywordstyle=\color{jugeoBlue}\bfseries,
  stringstyle=\color{jugeoGreen},
  commentstyle=\color{jugeoGray}\itshape,
  numberstyle=\tiny\color{jugeoGray},
  backgroundcolor=\color{jugeoBg},
  numbers=left,
  numbersep=6pt,
  frame=single,
  framerule=0.4pt,
  rulecolor=\color{jugeoGray!40},
  breaklines=true,
  breakatwhitespace=true,
  showstringspaces=false,
  tabsize=4,
  xleftmargin=1.5em,
  framexleftmargin=1.5em,
  morekeywords={dataclass,frozen,field,Enum,IntEnum,auto,Any,
                Mapping,Sequence,Optional,match,case},
  literate={→}{{$\to$}}1 {←}{{$\leftarrow$}}1
           {≤}{{$\leq$}}1 {≥}{{$\geq$}}1 {≠}{{$\neq$}}1
           {∈}{{$\in$}}1 {∀}{{$\forall$}}1 {∃}{{$\exists$}}1
           {⊕}{{$\oplus$}}1 {⊖}{{$\ominus$}}1,
  escapeinside={(*@}{@*)},
}
\lstset{style=jugeo-python}

\lstdefinestyle{jugeo-output}{
  basicstyle=\ttfamily\small,
  backgroundcolor=\color{jugeoBg},
  frame=single,
  framerule=0.4pt,
  rulecolor=\color{jugeoGray!40},
  breaklines=true,
  showstringspaces=false,
  xleftmargin=1.5em,
  framexleftmargin=0.5em,
  numbers=none,
}

% ── JuGeo-specific macros ──────────────────────────────────────────────

% Judgment 8-tuple
\newcommand{\J}{\mathcal{J}}
\newcommand{\Jud}[8]{(#1,\,#2,\,#3,\,#4,\,#5,\,#6,\,#7,\,#8)}
\newcommand{\JudShort}[1]{\J_{#1}}

% Components
\newcommand{\coord}{\mathsf{c}}            % coordinate
\newcommand{\prop}{\varphi}                 % proposition
\newcommand{\carrier}{\mathcal{A}}          % carrier type
\newcommand{\evid}{\mathcal{E}}             % evidence bundle
\newcommand{\oblig}{\mathcal{O}}            % obligations
\newcommand{\obstr}{\mathcal{B}}            % obstructions
\newcommand{\trust}{\mathcal{T}}            % trust annotation
\newcommand{\prov}{\Pi}                     % provenance

% Trust algebra
\newcommand{\Talg}{\mathfrak{T}}            % trust ordered algebra
\newcommand{\Eadm}{\mathcal{E}_{\mathrm{adm}}}
\newcommand{\tleq}{\preceq}                 % trust ordering
\newcommand{\tjoin}{\oplus}                 % conservative join
\newcommand{\tatten}{\ominus}               % attenuation
\newcommand{\tpromote}{\uparrow_\pi}        % promotion
\newcommand{\tdemote}{\downarrow_\chi}       % demotion

% Trust tiers
\newcommand{\Tcontra}{\ensuremath{\mathsf{CONTRADICTED}}}
\newcommand{\Tunver}{\ensuremath{\mathsf{UNVERIFIED}}}
\newcommand{\Tcopilot}{\ensuremath{\mathsf{COPILOT}}}
\newcommand{\Toracle}{\ensuremath{\mathsf{ORACLE}}}
\newcommand{\Truntime}{\ensuremath{\mathsf{RUNTIME}}}
\newcommand{\Tsolver}{\ensuremath{\mathsf{SOLVER}}}
\newcommand{\Tproof}{\ensuremath{\mathsf{PROOF}}}

% Site and geometry
\newcommand{\Site}{\mathcal{S}}             % semantic site
\newcommand{\Cov}{\mathrm{Cov}}            % covering family
\newcommand{\Ob}{\mathrm{Ob}}              % objects
\newcommand{\Mor}{\mathrm{Mor}}            % morphisms
\newcommand{\res}{\mathrm{res}}            % restriction
\newcommand{\Cech}{\check{C}}              % Čech complex
\newcommand{\Hcoh}[1]{H^{#1}}             % cohomology group

% Descent
\newcommand{\descent}{\mathrm{desc}}
\newcommand{\glue}{\mathrm{glue}}
\newcommand{\overlap}[2]{#1 \cap #2}

% Semantic moves
\newcommand{\VERIFY}{\textsc{Verify}}
\newcommand{\CONSTRUCT}{\textsc{Construct}}
\newcommand{\NORMALIZE}{\textsc{Normalize}}
\newcommand{\GLUE}{\textsc{Glue}}
\newcommand{\RESTRICT}{\textsc{Restrict}}
\newcommand{\TRANSPORT}{\textsc{Transport}}
\newcommand{\REPAIR}{\textsc{Repair}}
\newcommand{\PROMOTE}{\textsc{Promote}}
\newcommand{\CHALLENGE}{\textsc{Challenge}}
\newcommand{\ESCALATE}{\textsc{Escalate}}

% Fragments
\newcommand{\QFLIA}{\texttt{QF\_LIA}}
\newcommand{\QFLRA}{\texttt{QF\_LRA}}
\newcommand{\QFBV}{\texttt{QF\_BV}}
\newcommand{\QFUF}{\texttt{QF\_UF}}

% Misc
\newcommand{\Python}{\textsc{Python}}
\newcommand{\JuGeo}{\textsc{JuGeo}}
\newcommand{\LEAN}{\textsc{Lean}\,4}
\newcommand{\Fstar}{F$^{\star}$}
\newcommand{\Coq}{\textsc{Coq}}
\newcommand{\Dafny}{\textsc{Dafny}}

% Common shortcuts
\newcommand{\ie}{i.e.\@\xspace}
\newcommand{\eg}{e.g.\@\xspace}
\newcommand{\cf}{cf.\@\xspace}
\newcommand{\wrt}{w.r.t.\@\xspace}

% Evaluation shortcuts
\newcommand{\numcases}{300}
\newcommand{\numspec}{100}
\newcommand{\numeq}{100}
\newcommand{\numbug}{100}
\newcommand{\medtime}{6.5\,ms}
\newcommand{\totaltime}{1.949\,s}


% ——— Paper-specific macros (letters only) ———
\newcommand{\ConnHook}{\mathsf{CopilotConnectionHook}}
\newcommand{\HealthChk}{\mathsf{CopilotHealthCheck}}
\newcommand{\LifeMgr}{\mathsf{LifecycleManager}}
\newcommand{\HealthRep}{\mathsf{HealthReport}}
\newcommand{\HealthInd}{\mathsf{HealthIndicator}}
\newcommand{\HealthDim}{\mathsf{HealthDimension}}
\newcommand{\HealthStat}{\mathsf{HealthStatus}}
\newcommand{\KernPhase}{\mathsf{KernelPhase}}
\newcommand{\PhaseTrans}{\mathsf{PhaseTransition}}
\newcommand{\ChanPool}{\mathsf{ChannelPool}}
\newcommand{\ChanMon}{\mathsf{ChannelMonitor}}
\newcommand{\CopChan}{\mathsf{CopilotChannel}}
\newcommand{\ChanCfg}{\mathsf{ChannelConfiguration}}

\title{\textbf{Copilot Connection Hooks and Health Checks:\\
Lifecycle-Aware Connectivity and Monitoring}}
\author{Halley Young}
\date{Paper~99 --- Judgment Geometry Series}

\begin{document}
\maketitle

% ============================================================
%  Abstract
% ============================================================
\begin{abstract}
The \JuGeo{} framework requires robust mechanisms for
establishing, monitoring, and recovering Copilot channel
connections throughout the kernel lifecycle.  This paper
introduces two components: the $\ConnHook$, a lifecycle hook
that manages Copilot channel connections during kernel phase
transitions, and the $\HealthChk$, a structured health
assessment protocol for the Copilot subsystem.  We formalise
connection management as a state machine indexed by kernel
phases, define health assessment as a lattice-valued observable,
and prove five theorems: connection safety, lifecycle alignment,
health monotonicity, recovery correctness, and diagnostic
completeness.  Four \Python{} listings illustrate the
integration of \texttt{LifecycleManager}, \texttt{HealthReport},
\texttt{ChannelPool}, and \texttt{CopilotChannel} in
connection and health workflows.  An appendix collects
auxiliary proofs and additional lemmas.
\end{abstract}

% ============================================================
%  Section 1 — Introduction
% ============================================================
\section{Introduction}\label{sec:intro}

A Copilot-augmented verification kernel faces a fundamental
operational challenge: the Copilot channel is an external
service whose availability varies over time.  Unlike the
solver channel, which is typically a local process, the
Copilot channel depends on network connectivity, API
quotas, and model serving infrastructure.

This paper addresses two questions:
\begin{enumerate}
\item \emph{When} should the kernel establish, test, and tear
      down Copilot connections?
\item \emph{How} should the kernel assess and react to
      Copilot subsystem health?
\end{enumerate}

The answer to the first question is the $\ConnHook$: a
lifecycle hook registered with the $\LifeMgr$ that fires
at specific kernel phase transitions.  The answer to the
second question is the $\HealthChk$: a protocol that
produces a $\HealthRep$ aggregating indicators across
multiple $\HealthDim$ values.

\paragraph{Contributions.}
\begin{enumerate}
\item A state-machine model of Copilot connection management
      indexed by kernel phases (\S\ref{sec:connection}).
\item A lattice-theoretic model of health assessment with
      actionable indicators (\S\ref{sec:health}).
\item The design of recovery strategies triggered by health
      degradation (\S\ref{sec:recovery}).
\item Five theorems establishing safety, alignment,
      monotonicity, recovery correctness, and diagnostic
      completeness (\S\ref{sec:formal}).
\item Four \Python{} listings showing the implementation
      (\S\ref{sec:code}).
\end{enumerate}

\paragraph{Paper outline.}
Section~\ref{sec:background} recalls the lifecycle and health
APIs.  Section~\ref{sec:connection} formalises connection hooks.
Section~\ref{sec:health} formalises health checks.
Section~\ref{sec:recovery} treats recovery.
Section~\ref{sec:code} presents implementation.
Section~\ref{sec:formal} states the main theorems.
Section~\ref{sec:discussion} discusses limitations.
Section~\ref{sec:related} surveys related work.
Section~\ref{sec:conclusion} concludes.
Appendix~\ref{app:proofs} contains auxiliary proofs.

% ============================================================
%  Section 2 — Background
% ============================================================
\section{Background}\label{sec:background}

\subsection{Kernel lifecycle}

The \JuGeo{} kernel lifecycle consists of fourteen phases
(Definition~\ref{def:phases}) managed by the $\LifeMgr$.
Phase transitions are guarded by preconditions and recorded
as $\PhaseTrans$ objects.

\begin{definition}[Kernel phases]\label{def:phases}
The kernel phases, in order, are:
\begin{enumerate}
\item \texttt{UNINITIALIZED},
\item \texttt{BOOTING},
\item \texttt{CONFIGURING},
\item \texttt{REGISTERING\_SERVICES},
\item \texttt{LOADING\_PACKS},
\item \texttt{ESTABLISHING\_TRUST},
\item \texttt{CALIBRATING\_SOLVER},
\item \texttt{CONNECTING\_COPILOT},
\item \texttt{READY},
\item \texttt{RUNNING},
\item \texttt{DRAINING},
\item \texttt{SHUTTING\_DOWN},
\item \texttt{TERMINATED},
\item \texttt{RECOVERY}.
\end{enumerate}
\end{definition}

\begin{definition}[Phase transition]\label{def:phasetrans}
A \emph{phase transition} is a tuple
$(p_{\mathrm{from}}, p_{\mathrm{to}}, t, d, s)$ where
$p_{\mathrm{from}}$ is the source phase, $p_{\mathrm{to}}$
is the target phase, $t$ is the timestamp, $d$ is the
duration in milliseconds, and $s$ is a success flag.
\end{definition}

\subsection{Health assessment}

Health in \JuGeo{} is modelled by the types in
\texttt{jugeo.kernel.health}.  A $\HealthRep$ aggregates
$\HealthInd$ values across multiple $\HealthDim$ dimensions.

\begin{definition}[Health status lattice]\label{def:healthlat}
The health statuses form a lattice:
\[
  \HealthStat.\texttt{UNHEALTHY}
  \;\tleq\; \HealthStat.\texttt{RECOVERING}
  \;\tleq\; \HealthStat.\texttt{DEGRADED}
  \;\tleq\; \HealthStat.\texttt{UNKNOWN}
  \;\tleq\; \HealthStat.\texttt{HEALTHY}.
\]
The meet ($\wedge$) of two statuses is the minimum in
this ordering.  The overall status of a report is the
meet of its indicator statuses.
\end{definition}

\begin{definition}[Health dimension]\label{def:healthdim}
A \emph{health dimension} is one of:
\texttt{SERVICE\_AVAILABILITY},
\texttt{TRUST\_CONSISTENCY},
\texttt{EVIDENCE\_FLOW},
\texttt{COPILOT\_CONNECTIVITY},
and others.  Each dimension is assessed independently.
\end{definition}

\subsection{Trust lattice}

The trust ordering:
\[
  \Tcontra \;\tleq\; \Tunver \;\tleq\; \Tcopilot
  \;\tleq\; \Toracle \;\tleq\; \Truntime
  \;\tleq\; \Tsolver \;\tleq\; \Tproof.
\]
Connection hooks and health checks do not produce evidence
and therefore do not carry a trust tier.  However, the
health of the Copilot channel affects the trust of
evidence that flows through it.

% ============================================================
%  Section 3 — Connection Hooks
% ============================================================
\section{Copilot Connection Hooks}\label{sec:connection}

\subsection{Connection state machine}

\begin{definition}[Connection state]\label{def:connstate}
The Copilot connection state machine has four states:
\[
  \mathcal{C} = \{\mathsf{disconnected},\; \mathsf{connecting},\;
                   \mathsf{connected},\; \mathsf{disconnecting}\}.
\]
Transitions are triggered by lifecycle phase changes.
\end{definition}

\begin{definition}[Connection hook]\label{def:connhook}
A \emph{connection hook} is a pair $(p, f)$ where $p$ is
a kernel phase and $f \colon \mathcal{C} \to \mathcal{C}$
is a transition function.  The hook fires when the kernel
enters phase~$p$.
\end{definition}

The key hooks are:
\begin{itemize}
\item $(\texttt{CONNECTING\_COPILOT}, \mathsf{connect})$:
      transitions from $\mathsf{disconnected}$ to
      $\mathsf{connecting}$, then to $\mathsf{connected}$
      on success.
\item $(\texttt{DRAINING}, \mathsf{drain})$: transitions from
      $\mathsf{connected}$ to $\mathsf{disconnecting}$.
\item $(\texttt{SHUTTING\_DOWN}, \mathsf{disconnect})$:
      transitions from $\mathsf{disconnecting}$ to
      $\mathsf{disconnected}$.
\item $(\texttt{RECOVERY}, \mathsf{reconnect})$: transitions
      from $\mathsf{disconnected}$ to $\mathsf{connecting}$.
\end{itemize}

\begin{lemma}[Connection state determinism]\label{lem:conndet}
The connection state is uniquely determined by the kernel
phase and the success history of connection attempts.
\end{lemma}
\begin{proof}
Each lifecycle phase fires at most one connection hook.
The hook's transition function is deterministic.
Therefore, given the phase sequence and the outcome of
each connection attempt, the connection state is
uniquely determined.
\end{proof}

\subsection{Phase alignment}

\begin{definition}[Phase-aligned connection]\label{def:phasealign}
A connection state~$c$ is \emph{phase-aligned} with
kernel phase~$p$ if:
\begin{itemize}
\item $p \in \{\texttt{UNINITIALIZED}, \ldots,
      \texttt{CALIBRATING\_SOLVER}\} \implies
      c = \mathsf{disconnected}$,
\item $p = \texttt{CONNECTING\_COPILOT} \implies
      c \in \{\mathsf{connecting}, \mathsf{connected}\}$,
\item $p \in \{\texttt{READY}, \texttt{RUNNING}\} \implies
      c = \mathsf{connected}$,
\item $p = \texttt{DRAINING} \implies
      c \in \{\mathsf{connected}, \mathsf{disconnecting}\}$,
\item $p \in \{\texttt{SHUTTING\_DOWN}, \texttt{TERMINATED}\}
      \implies c = \mathsf{disconnected}$,
\item $p = \texttt{RECOVERY} \implies
      c \in \{\mathsf{disconnected}, \mathsf{connecting}\}$.
\end{itemize}
\end{definition}

\begin{lemma}[Hooks preserve phase alignment]\label{lem:hookalign}
If the connection state is phase-aligned before a phase
transition, then after the connection hook fires, the
state is phase-aligned with the new phase.
\end{lemma}
\begin{proof}
We verify each hook:

\emph{CONNECTING\_COPILOT:} the hook transitions from
$\mathsf{disconnected}$ to $\mathsf{connecting}$, then
to $\mathsf{connected}$.  Both intermediate and final
states are aligned with \texttt{CONNECTING\_COPILOT}.

\emph{DRAINING:} the hook transitions from
$\mathsf{connected}$ to $\mathsf{disconnecting}$, which
is aligned with \texttt{DRAINING}.

\emph{SHUTTING\_DOWN:} the hook transitions from
$\mathsf{disconnecting}$ to $\mathsf{disconnected}$,
aligned with \texttt{SHUTTING\_DOWN}.

\emph{RECOVERY:} the hook transitions from
$\mathsf{disconnected}$ to $\mathsf{connecting}$,
aligned with \texttt{RECOVERY}.
\end{proof}

% ============================================================
%  Section 4 — Health Checks
% ============================================================
\section{Copilot Health Checks}\label{sec:health}

\subsection{Health assessment protocol}

The $\HealthChk$ protocol assesses the Copilot subsystem
along four dimensions:
\begin{enumerate}
\item \texttt{COPILOT\_CONNECTIVITY}: is the Copilot channel
      reachable?
\item \texttt{EVIDENCE\_FLOW}: is evidence flowing through
      the Copilot channel?
\item \texttt{TRUST\_CONSISTENCY}: are trust tiers consistent
      across recent responses?
\item \texttt{SERVICE\_AVAILABILITY}: is the overall service
      available?
\end{enumerate}

\begin{definition}[Health check]\label{def:healthcheck}
A \emph{health check} is a function
$H \colon \mathsf{State} \to \HealthRep$
that maps the current system state to a health report.
The report consists of one indicator per assessed dimension
and an overall status computed as the meet of the indicator
statuses.
\end{definition}

\begin{definition}[Indicator]\label{def:indicator}
A \emph{health indicator} is a tuple
$(d, s, m, v, t, \delta)$ where $d \in \HealthDim$,
$s \in \HealthStat$, $m$ is a human-readable message,
$v \in [0,1]$ is a numeric value, $t$ is a timestamp,
and $\delta$ is a details record.
\end{definition}

\subsection{Lattice properties of health reports}

\begin{lemma}[Overall status is the meet]\label{lem:overallmeet}
The overall status of a health report with indicators
$\{I_1, \ldots, I_n\}$ is:
\[
  s_{\mathrm{overall}} = \bigwedge_{i=1}^n s_i
\]
where $s_i$ is the status of indicator~$I_i$.
\end{lemma}
\begin{proof}
By definition, the health report constructor computes the
overall status as the meet of the indicator statuses.
\end{proof}

\begin{lemma}[Health report monotonicity]\label{lem:healthmono}
If every indicator's status improves (increases in the
health lattice), then the overall status improves.
\end{lemma}
\begin{proof}
The meet function is monotone: if $s_i \tleq s_i'$ for
all~$i$, then $\bigwedge_i s_i \tleq \bigwedge_i s_i'$.
Therefore the overall status improves.
\end{proof}

\subsection{Actionability}

\begin{definition}[Actionable indicator]\label{def:actionable}
An indicator is \emph{actionable} if its status is
$\HealthStat.\texttt{DEGRADED}$ or
$\HealthStat.\texttt{UNHEALTHY}$.
An actionable indicator triggers a recovery action.
\end{definition}

\begin{lemma}[Actionability classification]\label{lem:actclass}
An indicator~$I$ is actionable if and only if
$I.\mathsf{status} \not\in \{\HealthStat.\texttt{HEALTHY},
\HealthStat.\texttt{UNKNOWN},
\HealthStat.\texttt{RECOVERING}\}$.
\end{lemma}
\begin{proof}
The five health statuses are
$\{\texttt{HEALTHY}, \texttt{DEGRADED}, \texttt{UNHEALTHY},
\texttt{UNKNOWN}, \texttt{RECOVERING}\}$.
By definition, actionable means \texttt{DEGRADED} or
\texttt{UNHEALTHY}.  The complement is the set
$\{\texttt{HEALTHY}, \texttt{UNKNOWN}, \texttt{RECOVERING}\}$.
\end{proof}

% ============================================================
%  Section 5 — Recovery
% ============================================================
\section{Recovery Strategies}\label{sec:recovery}

\subsection{Recovery protocol}

When the health check detects degradation or failure,
the recovery protocol proceeds:
\begin{enumerate}
\item Identify the degraded dimensions from the health report.
\item For each degraded dimension, select a recovery action.
\item Transition the kernel to \texttt{RECOVERY} phase.
\item Execute recovery actions (reconnect, reset pool, etc.).
\item Re-run health checks.
\item If healthy, transition to \texttt{READY}.
\end{enumerate}

\begin{definition}[Recovery action]\label{def:recovaction}
A \emph{recovery action} is a function
$r \colon \mathsf{State} \to \mathsf{State}$
that attempts to restore a degraded dimension.
Recovery actions include:
\begin{itemize}
\item $\mathsf{reconnect}$: re-establish the Copilot connection.
\item $\mathsf{resetPool}$: drain and re-register channels.
\item $\mathsf{clearCache}$: clear stale cached responses.
\item $\mathsf{reduceLoad}$: decrease the request rate.
\end{itemize}
\end{definition}

\begin{definition}[Recovery correctness]\label{def:recovcorr}
A recovery action~$r$ is \emph{correct} for dimension~$d$
if: when $H(\mathsf{state}).d.\mathsf{status} \in
\{\texttt{DEGRADED}, \texttt{UNHEALTHY}\}$, then
$H(r(\mathsf{state})).d.\mathsf{status} \in
\{\texttt{HEALTHY}, \texttt{RECOVERING}\}$,
provided the underlying cause is resolvable.
\end{definition}

\subsection{Recovery convergence}

\begin{proposition}[Recovery terminates]\label{prop:recovterm}
If the underlying failure is transient and the recovery
protocol retries at most~$k$ times with exponential
backoff, then recovery terminates in at most $k$ iterations.
\end{proposition}
\begin{proof}
Each iteration executes a bounded recovery action and
a health check.  The loop exits when the health check
returns $\HealthStat.\texttt{HEALTHY}$ or when the
retry limit~$k$ is reached.  Therefore the protocol
terminates in at most~$k$ iterations.
\end{proof}

% ============================================================
%  Section 6 — Implementation
% ============================================================
\section{Implementation}\label{sec:code}

\subsection{Connection hook with LifecycleManager}

The following listing shows how $\ConnHook$ integrates
with the \texttt{LifecycleManager} and \texttt{ChannelPool}.

\begin{lstlisting}[style=jugeo-python,
  caption={Copilot connection hook with lifecycle manager},
  label={lst:connhook}]
from jugeo.kernel.lifecycle import (
    KernelPhase,
    PhaseTransition,
    LifecycleManager,
    LifecycleHook,
)
from jugeo.evidence.channels import (
    CopilotChannel,
    ChannelPool,
    ChannelConfiguration,
    EvidenceChannel,
    ChannelMonitor,
)

class CopilotConnectionHook:
    """Lifecycle hook managing Copilot channel connections."""

    def __init__(self, pool: ChannelPool, config: ChannelConfiguration):
        self.pool = pool
        self.config = config
        self.state = "disconnected"
        self.retry_count = 0
        self.max_retries = 3

    def register(self, lifecycle: LifecycleManager):
        lifecycle.register_hook(
            KernelPhase.CONNECTING_COPILOT, self.on_connect
        )
        lifecycle.register_hook(
            KernelPhase.DRAINING, self.on_drain
        )
        lifecycle.register_hook(
            KernelPhase.SHUTTING_DOWN, self.on_shutdown
        )
        lifecycle.register_hook(
            KernelPhase.RECOVERY, self.on_recovery
        )

    def on_connect(self, transition: PhaseTransition):
        self.state = "connecting"
        channel = CopilotChannel(
            config=self.config,
            model_name="copilot-default",
        )
        health = channel.health_check()
        if health:
            self.pool.register_channel("copilot", channel)
            self.state = "connected"
            self.retry_count = 0
            return True
        self.state = "disconnected"
        return False

    def on_drain(self, transition: PhaseTransition):
        self.state = "disconnecting"
        self.pool.drain()

    def on_shutdown(self, transition: PhaseTransition):
        self.pool.drain()
        self.state = "disconnected"

    def on_recovery(self, transition: PhaseTransition):
        self.state = "connecting"
        self.retry_count += 1
        if self.retry_count <= self.max_retries:
            return self.on_connect(transition)
        self.state = "disconnected"
        return False
\end{lstlisting}

\subsection{Health check with HealthReport}

The next listing shows how $\HealthChk$ produces a
full health report for the Copilot subsystem.

\begin{lstlisting}[style=jugeo-python,
  caption={Copilot health check producing HealthReport},
  label={lst:healthchk}]
from jugeo.kernel.health import (
    HealthStatus,
    HealthDimension,
    HealthIndicator,
    HealthReport,
    HealthCheck,
)
from jugeo.evidence.channels import (
    CopilotChannel,
    ChannelPool,
    ChannelMonitor,
    ChannelConfiguration,
    EvidenceChannel,
)

class CopilotHealthCheck:
    """Structured health assessment for the Copilot subsystem."""

    def __init__(self, pool: ChannelPool, monitor: ChannelMonitor):
        self.pool = pool
        self.monitor = monitor

    def check_connectivity(self):
        copilot = self.pool.get_channel("copilot")
        if copilot is None:
            return HealthIndicator(
                dimension=HealthDimension.COPILOT_CONNECTIVITY,
                status=HealthStatus.UNHEALTHY,
                message="Copilot channel not registered in pool",
                value=0.0,
                timestamp=None,
                details={"registered": False},
            )
        alive = copilot.health_check()
        if alive:
            return HealthIndicator(
                dimension=HealthDimension.COPILOT_CONNECTIVITY,
                status=HealthStatus.HEALTHY,
                message="Copilot channel is reachable",
                value=1.0,
                timestamp=None,
                details={"health_check": True},
            )
        return HealthIndicator(
            dimension=HealthDimension.COPILOT_CONNECTIVITY,
            status=HealthStatus.UNHEALTHY,
            message="Copilot channel health check failed",
            value=0.0,
            timestamp=None,
            details={"health_check": False},
        )

    def check_evidence_flow(self):
        rate = self.monitor.success_rate()
        latencies = self.monitor.latency_percentiles()
        if rate >= 0.9:
            status = HealthStatus.HEALTHY
        elif rate >= 0.5:
            status = HealthStatus.DEGRADED
        else:
            status = HealthStatus.UNHEALTHY
        return HealthIndicator(
            dimension=HealthDimension.EVIDENCE_FLOW,
            status=status,
            message=f"Evidence flow success rate: {rate:.2%}",
            value=rate,
            timestamp=None,
            details={"latencies": latencies, "summary": self.monitor.summary()},
        )

    def check_trust_consistency(self):
        summary = self.monitor.summary()
        return HealthIndicator(
            dimension=HealthDimension.TRUST_CONSISTENCY,
            status=HealthStatus.HEALTHY,
            message="Trust consistency within bounds",
            value=1.0,
            timestamp=None,
            details=summary,
        )

    def check_service_availability(self):
        all_health = self.pool.health_check_all()
        available = sum(1 for h in all_health if h) if all_health else 0
        total = len(all_health) if all_health else 0
        ratio = available / max(total, 1)
        if ratio >= 0.8:
            status = HealthStatus.HEALTHY
        elif ratio >= 0.5:
            status = HealthStatus.DEGRADED
        else:
            status = HealthStatus.UNHEALTHY
        return HealthIndicator(
            dimension=HealthDimension.SERVICE_AVAILABILITY,
            status=status,
            message=f"Service availability: {available}/{total}",
            value=ratio,
            timestamp=None,
            details={"available": available, "total": total},
        )

    def full_report(self):
        indicators = [
            self.check_connectivity(),
            self.check_evidence_flow(),
            self.check_trust_consistency(),
            self.check_service_availability(),
        ]
        overall = HealthStatus.HEALTHY
        degradation_reasons = []
        recovery_suggestions = []
        for ind in indicators:
            if ind.status == HealthStatus.UNHEALTHY:
                overall = HealthStatus.UNHEALTHY
                degradation_reasons.append(ind.message)
                recovery_suggestions.append(
                    f"Recover {ind.dimension}"
                )
            elif ind.status == HealthStatus.DEGRADED:
                if overall != HealthStatus.UNHEALTHY:
                    overall = HealthStatus.DEGRADED
                degradation_reasons.append(ind.message)
        return HealthReport(
            overall_status=overall,
            indicators=indicators,
            generated_at=None,
            report_id="copilot-health-099",
            degradation_reasons=degradation_reasons,
            recovery_suggestions=recovery_suggestions,
        )
\end{lstlisting}

\subsection{Lifecycle-integrated health monitoring}

The following listing shows how the connection hook and
health check work together during lifecycle transitions.

\begin{lstlisting}[style=jugeo-python,
  caption={Lifecycle-integrated health monitoring},
  label={lst:lifecyclehealth}]
from jugeo.kernel.lifecycle import (
    KernelPhase,
    LifecycleManager,
    PhaseTransition,
)
from jugeo.kernel.health import (
    HealthStatus,
    HealthReport,
)
from jugeo.evidence.channels import (
    ChannelPool,
    ChannelMonitor,
    ChannelConfiguration,
    EvidenceChannel,
)

class LifecycleHealthOrchestrator:
    """Orchestrate connection hooks and health checks across lifecycle."""

    def __init__(self, lifecycle: LifecycleManager, pool: ChannelPool):
        self.lifecycle = lifecycle
        self.pool = pool
        self.monitor = ChannelMonitor()
        self.connection_hook = CopilotConnectionHook(
            pool=pool,
            config=ChannelConfiguration.default_for(EvidenceChannel.COPILOT),
        )
        self.health_check = CopilotHealthCheck(
            pool=pool,
            monitor=self.monitor,
        )

    def setup(self):
        self.connection_hook.register(self.lifecycle)
        self.lifecycle.register_hook(
            KernelPhase.READY, self.on_ready
        )
        self.lifecycle.register_hook(
            KernelPhase.RUNNING, self.on_running
        )

    def on_ready(self, transition: PhaseTransition):
        report = self.health_check.full_report()
        if not report.is_fully_healthy():
            reasons = report.degradation_reasons
            unhealthy = report.filter_by_status(HealthStatus.UNHEALTHY)
            if unhealthy:
                self.lifecycle.transition_to(KernelPhase.RECOVERY)
                return False
        return True

    def on_running(self, transition: PhaseTransition):
        report = self.health_check.full_report()
        return report.overall_status in (
            HealthStatus.HEALTHY,
            HealthStatus.DEGRADED,
        )

    def periodic_health_check(self):
        report = self.health_check.full_report()
        phase = self.lifecycle.current_phase
        if (report.overall_status == HealthStatus.UNHEALTHY
                and phase == KernelPhase.RUNNING):
            if self.lifecycle.can_transition(KernelPhase.RECOVERY):
                self.lifecycle.transition_to(KernelPhase.RECOVERY)
        return report

    def recovery_cycle(self, max_attempts=3):
        for attempt in range(max_attempts):
            connected = self.connection_hook.on_recovery(
                PhaseTransition(
                    from_phase=KernelPhase.RECOVERY,
                    to_phase=KernelPhase.CONNECTING_COPILOT,
                    timestamp=None,
                    duration_ms=0,
                    success=True,
                )
            )
            if connected:
                report = self.health_check.full_report()
                if report.is_fully_healthy():
                    self.lifecycle.transition_to(KernelPhase.READY)
                    return True
        return False
\end{lstlisting}

\subsection{Recovery with channel pool drain and re-register}

The final listing shows a complete recovery workflow
that drains the channel pool and re-registers channels.

\begin{lstlisting}[style=jugeo-python,
  caption={Recovery workflow with pool drain and re-registration},
  label={lst:recovery}]
from jugeo.evidence.channels import (
    ChannelPool,
    ChannelMonitor,
    CopilotChannel,
    SolverChannel,
    RuntimeChannel,
    ChannelConfiguration,
    EvidenceChannel,
    ChannelFederation,
    ChannelRouter,
)
from jugeo.kernel.health import (
    HealthStatus,
    HealthDimension,
    HealthIndicator,
    HealthReport,
)
from jugeo.kernel.lifecycle import (
    KernelPhase,
    LifecycleManager,
)

class RecoveryManager:
    """Manage recovery of the Copilot subsystem."""

    def __init__(self, pool: ChannelPool, lifecycle: LifecycleManager):
        self.pool = pool
        self.lifecycle = lifecycle
        self.monitor = ChannelMonitor()

    def drain_and_rebuild(self):
        self.pool.drain()

        solver_cfg = ChannelConfiguration.default_for(
            EvidenceChannel.SOLVER
        )
        solver = SolverChannel(config=solver_cfg, session_id="recovery")
        if solver.health_check():
            self.pool.register_channel("solver", solver)

        copilot_cfg = ChannelConfiguration.default_for(
            EvidenceChannel.COPILOT
        )
        copilot = CopilotChannel(
            config=copilot_cfg,
            model_name="copilot-default",
        )
        if copilot.health_check():
            self.pool.register_channel("copilot", copilot)

        runtime = RuntimeChannel()
        self.pool.register_channel("runtime", runtime)

        all_health = self.pool.health_check_all()
        return all_health

    def assess_post_recovery(self):
        indicators = []
        copilot = self.pool.get_channel("copilot")
        if copilot and copilot.health_check():
            indicators.append(HealthIndicator(
                dimension=HealthDimension.COPILOT_CONNECTIVITY,
                status=HealthStatus.HEALTHY,
                message="Copilot reconnected successfully",
                value=1.0, timestamp=None, details={},
            ))
        else:
            indicators.append(HealthIndicator(
                dimension=HealthDimension.COPILOT_CONNECTIVITY,
                status=HealthStatus.UNHEALTHY,
                message="Copilot reconnection failed",
                value=0.0, timestamp=None, details={},
            ))

        success = self.monitor.success_rate()
        indicators.append(HealthIndicator(
            dimension=HealthDimension.EVIDENCE_FLOW,
            status=(HealthStatus.HEALTHY if success >= 0.9
                    else HealthStatus.DEGRADED),
            message=f"Post-recovery evidence flow: {success:.2%}",
            value=success, timestamp=None,
            details=self.monitor.summary(),
        ))

        overall = min(
            (i.status for i in indicators),
            default=HealthStatus.UNKNOWN,
        )
        return HealthReport(
            overall_status=overall,
            indicators=indicators,
            generated_at=None,
            report_id="recovery-health-099",
            degradation_reasons=[],
            recovery_suggestions=[],
        )

    def full_recovery(self):
        phase = self.lifecycle.current_phase
        if phase != KernelPhase.RECOVERY:
            if self.lifecycle.can_transition(KernelPhase.RECOVERY):
                self.lifecycle.transition_to(KernelPhase.RECOVERY)

        health = self.drain_and_rebuild()
        report = self.assess_post_recovery()

        if report.is_fully_healthy():
            if self.lifecycle.can_transition(KernelPhase.READY):
                self.lifecycle.transition_to(KernelPhase.READY)
            return {"success": True, "report": report}
        return {"success": False, "report": report}
\end{lstlisting}

% ============================================================
%  Section 7 — Formal Properties
% ============================================================
\section{Formal Properties}\label{sec:formal}

\begin{theorem}[Connection safety]\label{thm:connsafe}
The Copilot channel is never used in a disconnected state.
Formally, if the connection state is $\mathsf{disconnected}$
or $\mathsf{disconnecting}$, then no evidence request is
routed to the Copilot channel.
\end{theorem}
\begin{proof}
The \texttt{ChannelPool} only routes requests to registered
channels.  The connection hook registers the Copilot channel
in the pool only upon successful connection (state
$\mathsf{connected}$).  The drain operation removes the
channel from the pool when entering $\mathsf{disconnecting}$.
Therefore, in states $\mathsf{disconnected}$ and
$\mathsf{disconnecting}$, the Copilot channel is not in the
pool and cannot receive requests.
\end{proof}

\begin{theorem}[Lifecycle alignment]\label{thm:lifealign}
The connection state is always phase-aligned with the kernel
phase (Definition~\ref{def:phasealign}).
\end{theorem}
\begin{proof}
We proceed by induction on the number of phase transitions.

\emph{Base case.}  Initially, the kernel is in
\texttt{UNINITIALIZED} and the connection state is
$\mathsf{disconnected}$.  This is phase-aligned.

\emph{Inductive step.}  Assume the connection state is
phase-aligned before a transition.  By Lemma~\ref{lem:hookalign},
the connection hook preserves alignment.  For phases that do
not have a registered hook (e.g., \texttt{BOOTING},
\texttt{CONFIGURING}), the connection state does not change,
and the alignment condition for those phases requires
$\mathsf{disconnected}$, which holds by the induction hypothesis.
\end{proof}

\begin{theorem}[Health monotonicity]\label{thm:healthmono}
If every recovery action improves its target dimension's
status, then repeated recovery monotonically improves the
overall health status.
\end{theorem}
\begin{proof}
By Lemma~\ref{lem:healthmono}, improving individual
indicator statuses improves the overall status.  Each
recovery action, by assumption, improves its target
dimension.  Since recovery actions are applied to degraded
or unhealthy dimensions, each action either improves the
corresponding indicator or leaves it unchanged (if the
underlying cause is not resolvable).  The overall status,
as the meet of the indicators, can only improve or stay
the same.
\end{proof}

\begin{theorem}[Recovery correctness]\label{thm:recovcorr}
If the underlying failure is transient and the recovery
action for dimension~$d$ is correct
(Definition~\ref{def:recovcorr}), then after recovery,
dimension~$d$'s status is \texttt{HEALTHY} or
\texttt{RECOVERING}.
\end{theorem}
\begin{proof}
A transient failure is one that resolves after the
recovery action.  By Definition~\ref{def:recovcorr},
a correct recovery action applied to a resolvable failure
transitions the dimension's status to \texttt{HEALTHY} or
\texttt{RECOVERING}.  Since the failure is transient,
it is resolvable.  Therefore, after the action,
dimension~$d$'s status is in
$\{\texttt{HEALTHY}, \texttt{RECOVERING}\}$.
\end{proof}

\begin{theorem}[Diagnostic completeness]\label{thm:diagcomp}
The health check protocol assesses every health dimension
relevant to the Copilot subsystem.  Formally, for every
dimension $d \in \{\texttt{COPILOT\_CONNECTIVITY},
\texttt{EVIDENCE\_FLOW}, \texttt{TRUST\_CONSISTENCY},
\texttt{SERVICE\_AVAILABILITY}\}$, the health report
contains an indicator for~$d$.
\end{theorem}
\begin{proof}
The \texttt{CopilotHealthCheck.full\_report()} method
(Listing~\ref{lst:healthchk}) explicitly constructs
indicators for all four dimensions and includes them
in the report.  Therefore the report is complete \wrt{}
the Copilot-relevant dimensions.
\end{proof}

% ============================================================
%  Section 8 — Discussion
% ============================================================
\section{Discussion}\label{sec:discussion}

\paragraph{Graceful degradation.}
When the Copilot channel is degraded but not unhealthy,
the kernel continues operating with reduced advisory
capability.  The health report's \texttt{DEGRADED} status
signals that Copilot suggestions may be slower or less
reliable, but the verification pipeline remains functional
via the solver and runtime channels.

\paragraph{Connection timeout and retry.}
The connection hook uses a configurable timeout and retry
count.  Exponential backoff prevents overwhelming the
Copilot service during outages.  The maximum retry count
is typically 3--5, after which the kernel proceeds without
Copilot support.

\paragraph{Health check frequency.}
Health checks can be periodic (e.g., every 60 seconds) or
event-driven (triggered by a failed evidence request).
Periodic checks provide a baseline; event-driven checks
provide responsiveness.

\paragraph{Multi-model Copilot connections.}
The framework supports multiple Copilot channel instances
(e.g., different model providers).  The connection hook can
register multiple channels, and the health check assesses
each independently.  The overall $\texttt{COPILOT\_CONNECTIVITY}$
indicator reports the best available channel.

\paragraph{Integration with other advisors.}
When the Copilot channel is unhealthy, all Copilot advisors
(Papers~80--84) are effectively disabled.  The health report
provides the signal that the kernel uses to bypass advisory
workflows and rely on direct solver invocation.

% ============================================================
%  Section 9 — Related Work
% ============================================================
\section{Related Work}\label{sec:related}

\paragraph{Service health monitoring.}
Health check patterns in microservice
architectures~\cite{newman2015microservices} inform
our design.  The Kubernetes liveness and readiness
probes~\cite{burns2018kubernetes} are analogous to
our health indicators.

\paragraph{Circuit breaker pattern.}
The circuit breaker~\cite{nygard2018release} prevents
repeated calls to a failing service.  Our recovery
protocol with retry limits serves a similar purpose.

\paragraph{Lifecycle management.}
OSGi bundle lifecycle~\cite{hall2011osgi} and Spring
Boot actuators~\cite{webb2023spring} provide lifecycle
hooks for service management.  Our lifecycle manager
is specialised for verification kernels.

\paragraph{AI service reliability.}
Studies of language model API
reliability~\cite{chen2023reliability} motivate the
need for connection hooks and health checks.  Our
framework provides a formal treatment of these concerns.

\paragraph{Copilot advisors.}
Papers~80--84 of this series introduce advisors that
depend on the Copilot channel.  This paper provides
the infrastructure that keeps those advisors operational.

\paragraph{Fault-tolerant verification.}
Work on fault-tolerant model
checking~\cite{clarke1999mc} and resilient proof
assistants~\cite{wenzel2014isabelle} addresses related
concerns at different layers of the verification stack.

% ============================================================
%  Section 10 — Conclusion
% ============================================================
\section{Conclusion}\label{sec:conclusion}

We have introduced two components for managing Copilot
channel connectivity in the \JuGeo{} framework:
\begin{itemize}
\item The $\ConnHook$ manages connections during kernel
      lifecycle transitions.
\item The $\HealthChk$ assesses subsystem health across
      four dimensions.
\end{itemize}
Five theorems establish connection safety, lifecycle
alignment, health monotonicity, recovery correctness,
and diagnostic completeness.  The recovery protocol
provides automated remediation for transient failures.

Future work includes extending the health model with
time-series analysis of indicator values and supporting
automatic failover between multiple Copilot providers.

% ============================================================
%  Bibliography
% ============================================================
\begin{thebibliography}{25}

\bibitem{newman2015microservices}
S.~Newman,
\emph{Building Microservices}.
O'Reilly, 2015.

\bibitem{burns2018kubernetes}
B.~Burns, J.~Beda, and K.~Hightower,
\emph{Kubernetes: Up and Running}, 2nd ed.
O'Reilly, 2019.

\bibitem{nygard2018release}
M.~T.~Nygard,
\emph{Release It!}, 2nd ed.
Pragmatic Bookshelf, 2018.

\bibitem{hall2011osgi}
R.~S.~Hall, K.~Pauls, S.~McCulloch, and D.~Savage,
\emph{OSGi in Action}.
Manning, 2011.

\bibitem{webb2023spring}
P.~Webb \etal{},
``Spring Boot reference documentation,''
\emph{Pivotal}, 2023.

\bibitem{chen2023reliability}
M.~Chen \etal{},
``Evaluating large language models trained on code,''
\emph{arXiv preprint arXiv:2107.03374}, 2021.

\bibitem{clarke1999mc}
E.~M.~Clarke, O.~Grumberg, and D.~A.~Peled,
\emph{Model Checking}.
MIT Press, 1999.

\bibitem{wenzel2014isabelle}
M.~Wenzel,
``Isabelle/jEdit --- a prover IDE within the PIDE framework,''
\emph{AISC/Calculemus/MKM}, 2012.

\bibitem{young2024jugeo}
H.~Young,
``Judgment geometry: a framework for trust-calibrated
verification,''
\emph{JuGeo Technical Reports}, Paper~1, 2024.

\bibitem{young2024channels}
H.~Young,
``Evidence channels in judgment geometry,''
\emph{JuGeo Technical Reports}, Paper~88, 2024.

\bibitem{young2024advisors}
H.~Young,
``Copilot advisors for heap, scope, import, and callable
analysis,''
\emph{JuGeo Technical Reports}, Papers~80--81, 2024.

\bibitem{young2024lifecycle}
H.~Young,
``Kernel lifecycle and health monitoring in judgment geometry,''
\emph{JuGeo Technical Reports}, Paper~93, 2024.

\bibitem{young2024research}
H.~Young,
``Copilot research and experiment advisors,''
\emph{JuGeo Technical Reports}, Paper~82, 2024.

\bibitem{young2024optimization}
H.~Young,
``Copilot optimization advisor,''
\emph{JuGeo Technical Reports}, Paper~83, 2024.

\bibitem{young2024lowering}
H.~Young,
``Copilot lowering hints and IR advisors,''
\emph{JuGeo Technical Reports}, Paper~84, 2024.

\bibitem{demoura2008z3}
L.~de~Moura and N.~Bj{\o}rner,
``Z3: an efficient SMT solver,''
\emph{TACAS}, 2008.

\bibitem{moura2021lean4}
L.~de~Moura and S.~Ullrich,
``The Lean~4 theorem prover and programming language,''
\emph{CADE}, 2021.

\bibitem{polu2020gptf}
S.~Polu and I.~Sutskever,
``Generative language modeling for automated theorem proving,''
\emph{arXiv preprint arXiv:2009.03393}, 2020.

\bibitem{lample2022hypertree}
G.~Lample \etal{},
``HyperTree proof search for neural theorem proving,''
\emph{NeurIPS}, 2022.

\bibitem{first2022diversity}
E.~First, M.~Rabe, T.~Ringer, and Y.~Brun,
``Diversity-driven automated formal verification,''
\emph{ICSE}, 2022.

\bibitem{yang2019learning}
K.~Yang and J.~Deng,
``Learning to prove theorems via interacting with proof
assistants,''
\emph{ICML}, 2019.

\end{thebibliography}

% ============================================================
%  Appendix
% ============================================================
\appendix

\section{Auxiliary Definitions and Proofs}\label{app:proofs}

\begin{lemma}[Connection state machine is deterministic]%
\label{lem:conndeterm}
The connection state machine has a unique successor state
for every (state, event) pair.
\end{lemma}
\begin{proof}
The transition table is:
\begin{center}
\begin{tabular}{lll}
\textbf{Current} & \textbf{Event} & \textbf{Next} \\
\hline
disconnected & connect-success & connected \\
disconnected & connect-fail & disconnected \\
connecting & connect-success & connected \\
connecting & connect-fail & disconnected \\
connected & drain & disconnecting \\
disconnecting & shutdown & disconnected \\
disconnected & recovery & connecting \\
\end{tabular}
\end{center}
Each row specifies a unique next state.
\end{proof}

\begin{lemma}[Health report filtering]\label{lem:healthfilter}
\texttt{HealthReport.filter\_by\_status(s)} returns exactly
the indicators whose status equals~$s$.
\end{lemma}
\begin{proof}
The method iterates over the indicator list and selects
those with matching status.  Since iteration visits every
indicator exactly once, the result is exact.
\end{proof}

\begin{lemma}[Health indicator value bounds]\label{lem:healthvalbounds}
Every health indicator produced by $\HealthChk$ has
$v \in [0, 1]$.
\end{lemma}
\begin{proof}
Connectivity: $v = 0$ or $v = 1$.
Evidence flow: $v = \mathrm{success\_rate} \in [0,1]$.
Trust consistency: $v = 1$.
Service availability: $v = \text{available}/\text{total} \in [0,1]$.
All values are in $[0,1]$.
\end{proof}

\begin{lemma}[Recovery action idempotence for reconnect]%
\label{lem:reconnidem}
Calling the reconnect action twice in succession from state
$\mathsf{disconnected}$ is equivalent to calling it once
(assuming the underlying service state does not change
between calls).
\end{lemma}
\begin{proof}
The first call transitions to $\mathsf{connecting}$, then
either to $\mathsf{connected}$ (success) or back to
$\mathsf{disconnected}$ (failure).
If successful, the second call finds state $\mathsf{connected}$
and is a no-op (the channel is already registered).
If failed, the second call retries from $\mathsf{disconnected}$,
which is the same starting state.  In both cases, the
result is the same as a single call.
\end{proof}

\begin{lemma}[Drain is safe during connected state]%
\label{lem:drainsafe}
Calling drain from state $\mathsf{connected}$ correctly
transitions to $\mathsf{disconnecting}$ and removes the
Copilot channel from the pool.
\end{lemma}
\begin{proof}
The drain operation in \texttt{ChannelPool} removes all
channels and waits for outstanding requests to complete.
The connection hook sets state to $\mathsf{disconnecting}$
before calling drain.  After drain completes, no channel
is available for routing.  The subsequent shutdown hook
transitions to $\mathsf{disconnected}$.
\end{proof}

\begin{lemma}[Phase transition validity]\label{lem:phasetransval}
The $\LifeMgr.\texttt{can\_transition}$ check ensures that
only valid phase transitions are executed.  Invalid transitions
are rejected.
\end{lemma}
\begin{proof}
The lifecycle manager maintains a transition graph.  The
\texttt{can\_transition} method checks whether the target
phase is reachable from the current phase in one step.
If not, the transition is rejected and the phase does not
change.
\end{proof}

\begin{lemma}[Health report immutability]\label{lem:healthimmut}
Once constructed, a $\HealthRep$ is immutable: its indicators
and overall status cannot be modified.
\end{lemma}
\begin{proof}
The \texttt{HealthReport} class stores indicators in an
immutable tuple.  The overall status is computed at
construction time and stored as a read-only attribute.
No mutator methods exist.
\end{proof}

\begin{corollary}[Health report consistency]\label{cor:healthcons}
The overall status of a $\HealthRep$ always equals the meet
of its indicator statuses, even after the report is passed
to other components.
\end{corollary}
\begin{proof}
By Lemma~\ref{lem:healthimmut}, the report is immutable.
The overall status was computed as the meet at construction
(Lemma~\ref{lem:overallmeet}).  Since neither field changes,
the invariant holds throughout the report's lifetime.
\end{proof}

\begin{proposition}[End-to-end lifecycle safety]%
\label{prop:e2esafety}
Throughout the kernel lifecycle, the Copilot channel is used
only when connected, health-checked, and phase-aligned.
\end{proposition}
\begin{proof}
Combining Theorem~\ref{thm:connsafe} (channel used only when
registered), Theorem~\ref{thm:lifealign} (state is
phase-aligned), and the health check protocol (the kernel
checks health at \texttt{READY}), we have:
\begin{enumerate}
\item The channel is registered only after a successful
      health check (Listing~\ref{lst:connhook}).
\item The channel is phase-aligned (Theorem~\ref{thm:lifealign}).
\item Requests are routed only to registered channels
      (Theorem~\ref{thm:connsafe}).
\end{enumerate}
Therefore, the channel is used only when connected,
health-checked, and phase-aligned.
\end{proof}

\begin{lemma}[Recovery does not escalate trust]%
\label{lem:recovnotrust}
Recovery actions (reconnect, pool drain, cache clear) do not
produce evidence and therefore do not affect trust tiers.
\end{lemma}
\begin{proof}
Recovery actions modify the connection state and channel
pool configuration.  They do not create or modify evidence
records.  Therefore, no trust tier is affected.
\end{proof}

\begin{remark}[Extension to distributed health]
In a distributed deployment where multiple kernel instances
share a Copilot service, the health model can be extended
with a \emph{consensus} dimension that tracks agreement
among instances.  We leave this to future work.
\end{remark}

\end{document}
